\addtocontents{toc}{\protect\Needspace{3\baselineskip}}
\originalchapter{2007 年正式期末考试}
B. Zwiebach, 2007 年 5 月 23 日, 限时三小时, 允许个人三页笔记, 原卷末两页为公式表. 四道主问题及全部子问实际公开, 以下为中文重编与独立解答. 末题的长列表按模式类型列全, 同模式费米指标取反对称, 玻色指标取对称.

\src{8.251 Final Q1, (a)--(d), 25 分.}
\begin{exercise}
量子开玻色弦在 $x^2,x^3$ 平面的厄米角动量为
\[
 J=-\ii\sum_{n>0}\frac1n
 (\alpha_{-n}^{(2)}\alpha_n^{(3)}-\alpha_{-n}^{(3)}\alpha_n^{(2)}).
\]
有 $1+\alpha'M^2=N=N_{23}+N'$, 其中
$N_{23}=\sum_{n>0}(\alpha_{-n}^{(2)}\alpha_n^{(2)}+
\alpha_{-n}^{(3)}\alpha_n^{(3)})$, $N'$ 为其余横向方向.
定 $\alpha_n=(\alpha_n^{(2)}+\ii\alpha_n^{(3)})/\sqrt2$,
$\bar\alpha_n=(\alpha_n^{(2)}-\ii\alpha_n^{(3)})/\sqrt2$, 所以 $\alpha_n^\dagger=\bar\alpha_{-n}$.
(a) 求新振子对易式并改写 $J,N_{23}$.
(b) 对有限激发态
$\prod_{k>0}(\alpha_{-k})^{\lambda_k}(\bar\alpha_{-k})^{\bar\lambda_k}$ 乘其他方向激发和动量真空, 给 $J,N_{23}$ 本征值.
(c) 建立并证明 $J$ 与 $1+\alpha'M^2$ 的不等式.
(d) 对 $(\alpha_{-1})^N|p\rangle$ 求两者. 原文称指数为正整数, 有限一般态需允许未激发模式的指数为零.
\end{exercise}
\begin{solution}
由两实方向独立振子代数,
$[\alpha_m,\alpha_n]=[\bar\alpha_m,\bar\alpha_n]=0$,
$[\alpha_m,\bar\alpha_n]=m\delta_{m+n,0}$. 反代定义得
\[
 J=\sum_{n>0}\frac{\alpha_{-n}\bar\alpha_n-\bar\alpha_{-n}\alpha_n}{n},
 \qquad N_{23}=\sum_{n>0}(\alpha_{-n}\bar\alpha_n+\bar\alpha_{-n}\alpha_n).
\]
$\alpha_{-k}$ 增加 $J$ 一单位, $\bar\alpha_{-k}$ 减少一单位, 两者都增加层数 $k$. 因而
\[
 J=\sum_k(\lambda_k-\bar\lambda_k),\quad
 N_{23}=\sum_k k(\lambda_k+\bar\lambda_k).
\]
对任何这些本征态,
$|J|\leq\sum_k(\lambda_k+\bar\lambda_k)\leq N_{23}\leq N=1+\alpha'M^2$.
第一步是三角不等式, 第二步用 $k\geq1$, 第三步用 $N'\geq0$. 饱和须全部激发同一个角动量号、均为第一模式且其他方向不激发.
故 $(\alpha_{-1})^N|p\rangle$ 有 $J=N=1+\alpha'M^2$, 即最大正角动量轨迹 $J=1+\alpha'M^2$. 负号极化也可饱和绝对值. 经典式没有截距 1, 在大层数极限才忽略这个有限量.
\end{solution}

\src{8.251 Final Q2, (a)--(d), 25 分.}
\begin{exercise}
四维 de Sitter 平坦切片度规为
$g_{00}=-1,g_{ij}=\ee^{2Ht}\delta_{ij}$, $c=\hbar=1$, $H>0$.
Nambu--Goto 作用量 $S=-T_s\int\sqrt{(\dot X\cdot X')^2-\dot X^2X'^2}\dd\tau\dd\sigma$, $T_s=1/(2\pi\alpha')$, 内积用此曲度规.
(a) 选 $X^0=t=\tau$, 用 $\boldsymbol X$ 的导数写作用量.
(b) 取圆弦 $\boldsymbol X=(r(t)\cos\sigma,r(t)\sin\sigma,0)$, $0\leq\sigma<2\pi$, 积分 $\sigma$ 给 $L(r,\dot r;t)$, 再用物理半径 $R=\ee^{Ht}r$ 改写.
(c) 对常 $R$ 求势 $V(R)$、绘图、给定义域、临界半径及稳定性.
(d) 用 $R,\dot R$ 表示哈密顿量并判断物理运动中是否守恒.
\end{exercise}
\begin{solution}
记 $a(t)=\ee^{Ht}$, 有
$h_{tt}=-1+a^2\dot{\boldsymbol X}^2$,
$h_{\sigma\sigma}=a^2\boldsymbol X'^2$,
$h_{t\sigma}=a^2\dot{\boldsymbol X}\cdot\boldsymbol X'$. 所以
\[
 S=-T_s\int\dd t\dd\sigma\,
 a|\boldsymbol X'|
 \sqrt{1-a^2\left[\dot{\boldsymbol X}^2-
 \frac{(\dot{\boldsymbol X}\cdot\boldsymbol X')^2}{\boldsymbol X'^2}\right]}.
\]
圆弦径向速度与切线垂直, $|\boldsymbol X'|=r$, 故
\[
 L=-2\pi T_s a r\sqrt{1-a^2\dot r^2}
 =-\frac{R}{\alpha'}\sqrt{1-(\dot R-HR)^2}.
\]
令 $K=1/\alpha'$、$v=\dot R-HR$. 常 $R$ 的势为
$V=KR\sqrt{1-H^2R^2}$, $0\leq R\leq H^{-1}$, 正则类时内点取严格不等号. 导数
$V'=K(1-2H^2R^2)/\sqrt{1-H^2R^2}$, 临界点 $R_*=1/(\sqrt2H)$ 是最大值. $V''(R_*)=-4KH$, 而动力学速度 Hessian $L_{\dot R\dot R}=KR/(1-v^2)^{3/2}>0$, 在平衡点为 $2K/H$.
所以线性扰动满足 $\delta\ddot R=2H^2\delta R$, 有增长解, 不是稳定静弦.

\begin{center}
\begin{tikzpicture}
\begin{axis}[width=.68\textwidth,height=4.2cm,xmin=0,xmax=1,ymin=0,ymax=.55,
 xlabel={$HR$},ylabel={$HV/K$},domain=0:1,samples=121]
 \addplot[very thick,structurecolor]{x*sqrt(1-x*x)};
 \addplot[only marks,mark=*]coordinates{(.70710678,.5)};
\end{axis}
\end{tikzpicture}
\end{center}

共轭动量为 $P_R=KRv/\sqrt{1-v^2}$, 因而
\[
 \mathcal H=P_R\dot R-L
 =\frac{KR[1+HR(\dot R-HR)]}{\sqrt{1-(\dot R-HR)^2}}.
\]
$L(R,\dot R)$ 无显式时间, 所以满足欧拉--拉格朗日方程的运动有 $\dd\mathcal H/\dd t=0$. 这不是原共动变量中显式含 $a(t)$ 的瞬时能量, 而是变到 $R$ 后相应 de Sitter 时间平移生成元. 在恒 $R$ 平衡处 $\mathcal H=V(R_*)$ 可再次核对.
\end{solution}

\src{8.251 Final Q3, 三小问合计 20 分.}
\begin{exercise}
(a) 将 $L_{-2}^\perp L_{-2}^\perp|0\rangle$ 化为正规序振子态, $|0\rangle$ 的横向动量为零.
(b) $N$ 张重合 D$p$ 膜在 $d$ 个空间维度有什么无质量场? 加重合 O$p$ 后如何改变? 原卷公式表规定 $\Omega_p$ 在真空交换端点标签.
(c) 临界玻色理论中两个 D2 膜直角相交, 一个沿 $x^1,x^2$, 另一个沿 $x^1,x^3$. 找两个膜间弦全部快子质量平方, 并说明它们所在的方向.
\end{exercise}
\begin{solution}
(a) $L_{-2}|0\rangle=\frac12\alpha_{-1}\cdot\alpha_{-1}|0\rangle$,
$[L_{-2},\alpha_{-1}^I]=\alpha_{-3}^I$, 故
\[
 L_{-2}^2|0\rangle=
 \left[\alpha_{-3}\cdot\alpha_{-1}
 +\tfrac14(\alpha_{-1}\cdot\alpha_{-1})^2\right]|0\rangle.
\]
两种类型均在层 4, 零模项因 $p_T=0$ 消失.

(b) 无投影时为 $U(N)$ 邻接表示的向量及 $d-p$ 个法向标量, 标签空间维数 $N^2$. 按原卷的 $\Omega_p$ 转置标签、正真空与振子层 1 的负号, 极化矩阵须反对称, 得 $SO(N)$ 规范场及同表示的标量, 标签数 $N(N-1)/2$. 其他 orientifold 符号与辛型标签会给不同群; 本题采用已给投影, 不宣称所有 O 平面均同一选择.

(c) 光锥在共同 $t,x^1$ 方向, $x^2,x^3$ 各为一个 ND 方向, 其余 22 个横向坐标为 DD, 两膜在交线处无伸长. 零点常数 $-1+2/16=-7/8$, 故
$\alpha'M^2=N-7/8$.
层 $N=0$ 有一个快子, $M^2=-7/(8\alpha')$;
层 $N=1/2$ 有两个半整数第一模式态, 各有 $M^2=-3/(8\alpha')$.
$N\geq1$ 已非快子. 这些为一个有向端点扇区的计数, 反向弦给共轭标签.
它们只有共同 $x^1$ 的连续空间动量, 为沿交线传播的 $1+1$ 维场; ND 方向无零模动量, 不能当作整个 D2 平面传播的自由粒子.
\end{solution}

\src{8.251 Final Q4, (a)--(d), 30 分.}
\begin{exercise}
异质 $SO(32)$ 闭弦右边是八横向超弦, 有八个横向指标 $I=2,\ldots,9$ 的玻色 $\alpha_{-n}^I$, NS 半整数 $b_{-r}^I$ 与 R 整数 $d_{-n}^I$, 标准投影为 NS$+$ 与 R$-$. 左边有八个玻色 $\bar\alpha_{-n}^I$ 及 32 个实费米子 $\lambda^A$, $A=1,\ldots,32$, 分为 NS$'$ 半整数与 R$'$ 整数. NS$'$ 真空规定 $( -1)^{F_L}=+1$; R$'$ 零模正负奇偶各占一半. 左边朴素对称序质量式为
\[
 \alpha'M_L^2=\tfrac12\sum_{n\ne0}\bar\alpha_{-n}^I\bar\alpha_n^I
 +\tfrac12\sum_r r\lambda_{-r}^A\lambda_r^A,
\]
R$'$ 取整数 $r=n$, NS$'$ 取半整数. 不写动量标签.
(a) 求 NS$'$ 正规序常数, 只保留偶费米激发, 列全前三个质量层及简并数.
(b) 求 R$'$ 正规序常数和正奇偶真空数, 在总正奇偶投影下列前两层.
(c) 异质弦物理 $\alpha'M^2=4k$ 需左右质量贡献同为 $k$. 左边与右 NS$+$ 张量积给玻色, 与右 R$-$ 给费米. 有无快子? 列无质量态与玻色场, 给 $\alpha'M^2=4$ 时玻色加费米总数.
(d) 求完整投影左边生成函数 $f_L(q)=\sum_r d(r)q^r$, 指数 $r=\alpha'M_L^2$.
\end{exercise}
\begin{solution}
(a) 八玻色给 $-8/24$, 32 个 NS 费米给 $-32/48$, 共 $-1$. 所以
$\alpha'M_L^2=N_B+N_F-1$. 偶数费米投影保留下表全部类型:
\begin{center}
\begin{tabular}{c l r}\toprule
$k$ & 模式类型 & 态数\\\midrule
$-1$ & $|\mathrm{NS}'\rangle$ &1\\
0 & $\bar\alpha_{-1}^I$ &8\\
0 & $\lambda_{-1/2}^{[A}\lambda_{-1/2}^{B]}$ &496\\
1 & $\bar\alpha_{-2}^I$ &8\\
1 & $\bar\alpha_{-1}^{(I}\bar\alpha_{-1}^{J)}$ &36\\
1 & $\bar\alpha_{-1}^I\lambda_{-1/2}^{[A}\lambda_{-1/2}^{B]}$ &3968\\
1 & $\lambda_{-3/2}^A\lambda_{-1/2}^B$ &1024\\
1 & $\lambda_{-1/2}^{[A}\lambda_{-1/2}^B\lambda_{-1/2}^C\lambda_{-1/2}^{D]}$ &35960\\\bottomrule
\end{tabular}
\end{center}
每行均作用于 NS$'$ 真空. 相同半整数模式的费米指标取严格不同子集; 不同模式的两个指标没有反对称减半. 三层总数为 1、504、40996.

(b) R$'$ 费米给 $+32/24$, 与玻色合计 $+1$, 所以
$\alpha'M_L^2=N_B+N_F+1$.
32 个实零模给 $2^{16}$ 维 Clifford 真空空间, 正奇偶为 $2^{15}=32768$ 态 $|s_+\rangle$. 最低 $k=1$ 只有这些态.
下一层 $k=2$ 有 $\bar\alpha_{-1}^I|s_+\rangle$ 和
$\lambda_{-1}^A|s_-\rangle$, 分别 $8\cdot32768$ 与 $32\cdot32768$, 总数 1310720. 奇费米激发必须配负奇偶零模, 保持总正, 不能删掉后一类.

(c) 右 NS$+$ 与 R$-$ 均无负质量层, 所以左 $k=-1$ 虽有态, 无法匹配, 物理异质弦无此快子.
无质量左边为八个 $\bar\alpha_{-1}^I$ 与 496 个 $\lambda_{-1/2}^{[A}\lambda_{-1/2}^{B]}$. 右 NS$+$ 为 $b_{-1/2}^J|\mathrm{NS}\rangle$, 右 R$-$ 为八个旋量零模态. 前一类张量积给 $8\times8=64$ 个普适玻色极化, 即引力子 35、二形式 28、膨胀子 1; 第二类给 496 个规范向量、各八极化, 是 $SO(32)$ 邻接表示. 同样 504 个左态与右旋量给 4032 个费米态, 分成引力超多重态及 gaugino 态.
在 $k=1$, 左总数为 $40996+32768=73764$, 右 NS$+$、R$-$ 各有 128. 故玻色与费米各 9441792, 合计
\[
 73764(128+128)=18883584.
\]
不能把左 R$'$ 最低层遗漏, 也不能将匹配条件换成独立左右任意组合.

(d) 每类模式因子相乘, 偶费米投影用正负两个乘积的和,
\begin{align*}
 f_L(q)=\prod_{n\geq1}(1-q^n)^{-8}\Bigg\{
 &\frac1{2q}\left[\prod_{n\geq1}(1+q^{n-1/2})^{32}
 +\prod_{n\geq1}(1-q^{n-1/2})^{32}\right]\\
 &+32768q\prod_{n\geq1}(1+q^n)^{32}\Bigg\}.
\end{align*}
R 非零费米激发的奇偶只改变所配零模手征, 每种配置始终有 32768 个允许零模, 因此不需要再删掉半数正频配置. 有限展开给 $q^{-1}+504+73764q+\cdots$, 与逐类枚举独立对照.
\end{solution}

